\documentclass[aspectratio=169,11pt]{beamer}
\usepackage{../phanes-beamer}
\renewcommand{\ModuleID}{02}
\title{Responsible nuclear AI}
\subtitle{Specific regulatory questions, information flows, and review evidence}
\author{PHANES Initiative}
\institute{The University of Texas at Austin}
\date{September 12, 2026}
\begin{document}
\begin{frame}[plain]
\titlepage
\end{frame}
\begin{frame}[t,fragile]{Start with the actual transfer, not the product label}
\fontsize{11.5}{14}\selectfont
\begin{itemize}\setlength{\itemsep}{2mm}
\item Specify the information and activity: published report, input deck, geometry, design analysis, debugging assistance, or derived result.
\item Identify the people and organizations that may receive or access it, including support and administrators.
\item Map the endpoint, tools, logs, storage, backups, and onward sharing.
\item Determine the applicable regulatory route and project restrictions before using a real nonpublic case.
\end{itemize}
\vfill{\fontsize{6.5}{8}\selectfont\color{UTgray} Sources: utexport. See references and instructor notes.\par}
\end{frame}
\begin{frame}[t,fragile]{Three federal routes are not interchangeable}
\fontsize{10.5}{12.5}\selectfont\setlength{\tabcolsep}{4pt}\renewcommand{\arraystretch}{1.15}
\begin{tabular}{>{\raggedright\arraybackslash}p{\dimexpr 0.23\linewidth-3.68pt\relax}>{\raggedright\arraybackslash}p{\dimexpr 0.77\linewidth-12.32pt\relax}}
\toprule
\textbf{Route} & \textbf{Specific screening question} \\ \midrule
DOE / Part 810 & Is the activity covered nuclear assistance or a technology transfer within the rule's scope? \\[2pt]
NRC / Part 110 & Does the transaction involve covered nuclear equipment or material exports/imports? \\[2pt]
Commerce / EAR & Is the relevant software, technology, equipment, or release subject to Commerce jurisdiction? \\[2pt]
Project restrictions & Do contracts, confidentiality, sponsor conditions or institutional policies add limits? \\[2pt]
\bottomrule\end{tabular}\par\vspace{2mm}
\vfill{\fontsize{6.5}{8}\selectfont\color{UTgray} Sources: 810scope, 110, earrelease. See references and instructor notes.\par}
\end{frame}
\begin{frame}[t,fragile]{Part 810 scope includes reactor-development technology}
\fontsize{11.5}{14}\selectfont
\begin{itemize}\setlength{\itemsep}{2mm}
\item Section 810.2 addresses covered participation in foreign special-nuclear-material activities and related technology transfers.
\item Its listed activities include reactor development and specified reactor-related components, as well as fuel-cycle work.
\item Relevant technology can move through design analysis, instruction, consulting, or technical data.
\item A model-assisted review of a nonpublic reactor input deck is therefore a question to screen, even if the API server is physically in the United States.
\end{itemize}
\vfill{\fontsize{6.5}{8}\selectfont\color{UTgray} Sources: 810scope, 810definitions. See references and instructor notes.\par}
\end{frame}
\begin{frame}[t,fragile]{Read the rule in a useful order}
\fontsize{10.5}{12.5}\selectfont\setlength{\tabcolsep}{4pt}\renewcommand{\arraystretch}{1.15}
\begin{tabular}{>{\raggedright\arraybackslash}p{\dimexpr 0.2\linewidth-3.2pt\relax}>{\raggedright\arraybackslash}p{\dimexpr 0.8\linewidth-12.8pt\relax}}
\toprule
\textbf{Provision} & \textbf{What it helps establish} \\ \midrule
810.2 & Scope and exclusions; identify whether the material/activity is covered. \\[2pt]
810.3 & Definitions that give legal meaning to public availability, research and assistance. \\[2pt]
810.6 / 810.7 & General-authorization conditions versus activities requiring specific authorization. \\[2pt]
810.12 & Reporting obligations associated with the applicable authorized activity. \\[2pt]
810.5 & How to request DOE advice or an authoritative written interpretation. \\[2pt]
\bottomrule\end{tabular}\par\vspace{2mm}
\vfill{\fontsize{6.5}{8}\selectfont\color{UTgray} Sources: 810scope, 810definitions, 810general, 810specific, 810reports, 810interpret. See references and instructor notes.\par}
\end{frame}
\begin{frame}[t,fragile]{Screening an engineering workflow under Part 810}
\fontsize{10.5}{12.5}\selectfont\setlength{\tabcolsep}{4pt}\renewcommand{\arraystretch}{1.15}
\begin{tabular}{>{\raggedright\arraybackslash}p{\dimexpr 0.22\linewidth-3.52pt\relax}>{\raggedright\arraybackslash}p{\dimexpr 0.78\linewidth-12.48pt\relax}}
\toprule
\textbf{Screening step} & \textbf{Assessment and routing criterion} \\ \midrule
1. Scope screening & Does the workflow touch nuclear fuel-cycle or reactor technology under \S{}810.2? If not, Part 810 does not apply. \\[2pt]
2. Exclusion check & Is the material publicly available with provenance, or fundamental research under \S{}810.2(c)(2) and \S{}810.3? If yes, it is excluded. \\[2pt]
3. General authorization & Is the activity listed under \S{}810.6 (Appendix A destination or NRC licensee employee access)? If yes, general authorization applies. \\[2pt]
4. Reporting check & Generally authorized activities require a \S{}810.12 report within 30 days of initiation. \\[2pt]
5. Specific authorization & Non-Appendix destinations or sensitive technologies require DOE specific authorization under \S{}810.7 before any transfer. \\[2pt]
\bottomrule\end{tabular}\par\vspace{2mm}
\vfill{\fontsize{6.5}{8}\selectfont\color{UTgray} Sources: 810scope, 810general, 810specific, 810reports. See references and instructor notes.\par}
\end{frame}
\begin{frame}[t,fragile]{The public-information exclusion needs provenance}
\fontsize{11.5}{14}\selectfont
\begin{itemize}\setlength{\itemsep}{2mm}
\item Section 810.2(c)(2) excludes certain public information/technology transfers and fundamental-research results from Part 810.
\item Retain the publication, source URL, date, and evidence that the exact material is publicly available.
\item Do not assume that private annotations, unpublished dimensions, or partner results inherit the status of the public source.
\item For a mixed prompt, inventory each component and remove or review the nonpublic additions.
\end{itemize}
\vfill{\fontsize{6.5}{8}\selectfont\color{UTgray} Sources: 810scope, 810definitions. See references and instructor notes.\par}
\end{frame}
\begin{frame}[t,fragile]{Fundamental research is a defined condition}
\fontsize{11.5}{14}\selectfont
\begin{itemize}\setlength{\itemsep}{2mm}
\item Part 810's definition concerns research whose results are ordinarily published and broadly shared, distinguished from restricted proprietary work.
\item "At a university" and "we plan to publish later" do not resolve every input, contract, or disclosure restriction.
\item For EAR analysis, \S{}734.8 also distinguishes research outputs from controlled material supplied to conduct the research.
\item Bring the sponsor agreement and publication restrictions to the review; do not classify a project solely from its academic title.
\end{itemize}
\vfill{\fontsize{6.5}{8}\selectfont\color{UTgray} Sources: 810definitions, earresearch. See references and instructor notes.\par}
\end{frame}
\begin{frame}[t,fragile]{Worked case: public paper plus partner dimensions}
\fontsize{11.5}{14}\selectfont
\begin{itemize}\setlength{\itemsep}{2mm}
\item Input A: a public reactor-physics paper. Input B: a partner's unpublished fuel-layout table under an agreement.
\item The intended prompt asks the API to compare the partner design with the paper.
\item The transfer package now includes B and any excerpts derived from it; a citation to A does not establish B's status.
\item Proposed interim action: use a synthetic layout for the class example and send the mixed-case facts for review.
\end{itemize}
\end{frame}
\begin{frame}[t,fragile]{General authorization still has conditions}
\fontsize{11.5}{14}\selectfont
\begin{itemize}\setlength{\itemsep}{2mm}
\item Section 810.6 provides general authorization for specified activities when its conditions and exclusions are satisfied.
\item The Appendix country/entity route is subject to the rule's limitations, including the sensitive-technology restrictions.
\item The NRC-facility employment/access route has its own explicit conditions; it is not a blanket rule for all university collaborations.
\item Document which paragraph applies and check associated reporting and onward-transfer obligations.
\end{itemize}
\vfill{\fontsize{6.5}{8}\selectfont\color{UTgray} Sources: 810general. See references and instructor notes.\par}
\end{frame}
\begin{frame}[t,fragile]{Specific authorization must precede covered activity}
\fontsize{11.5}{14}\selectfont
\begin{itemize}\setlength{\itemsep}{2mm}
\item Section 810.7 identifies circumstances requiring a specific DOE authorization before the activity.
\item These include certain non-Appendix destinations and sensitive technologies, plus the activities enumerated in that section.
\item A vendor's general claim of US hosting does not identify the applicable authorization route.
\item If the route is unresolved, prepare the facts and use an approved public/synthetic alternative while the institution resolves it.
\end{itemize}
\vfill{\fontsize{6.5}{8}\selectfont\color{UTgray} Sources: 810specific. See references and instructor notes.\par}
\end{frame}
\begin{frame}[t,fragile]{Authorization and reporting are separate checks}
\fontsize{11.5}{14}\selectfont
\begin{itemize}\setlength{\itemsep}{2mm}
\item Section 810.12 includes reports due within 30 calendar days after beginning generally or specifically authorized activities, subject to its provisions.
\item The section also addresses completion/termination and other reporting situations.
\item Assign the reporting owner, start-date record, recipient details, and evidence of any required agreement.
\item Students document the workflow; the institution determines who submits and maintains the official record.
\end{itemize}
\vfill{\fontsize{6.5}{8}\selectfont\color{UTgray} Sources: 810reports. See references and instructor notes.\par}
\end{frame}
\begin{frame}[t,fragile]{US location is one fact among several}
\fontsize{10.5}{12.5}\selectfont\setlength{\tabcolsep}{4pt}\renewcommand{\arraystretch}{1.15}
\begin{tabular}{>{\raggedright\arraybackslash}p{\dimexpr 0.27\linewidth-4.32pt\relax}>{\raggedright\arraybackslash}p{\dimexpr 0.73\linewidth-11.68pt\relax}}
\toprule
\textbf{Question} & \textbf{Why it changes the review} \\ \midrule
Where is processing? & Identifies the computation and possible cross-border transmission. \\[2pt]
Who can access it? & Administrative, support, collaborator and subcontractor access may matter. \\[2pt]
What can be recovered? & Prompt logs, snapshots, backups and crash dumps can contain the same material. \\[2pt]
Where can it go next? & Agent tools, replication, export features and onward sharing create additional paths. \\[2pt]
\bottomrule\end{tabular}\par\vspace{2mm}
\end{frame}
\begin{frame}[t,fragile]{A domestic disclosure can still require analysis}
\fontsize{11.5}{14}\selectfont
\begin{itemize}\setlength{\itemsep}{2mm}
\item EAR \S{}734.13 includes certain releases of technology or source code to a foreign person inside the United States in its export definition.
\item Part 810 has its own scope, recipient definitions and exclusions; do not import EAR shorthand as the complete Part 810 test.
\item Nationality alone is not a sufficient classroom decision rule. The material, activity, and legal status also matter.
\item Have the institution assess the actual team and access arrangement rather than asking an agent to decide eligibility.
\end{itemize}
\vfill{\fontsize{6.5}{8}\selectfont\color{UTgray} Sources: earrelease, 810scope. See references and instructor notes.\par}
\end{frame}
\begin{frame}[t,fragile]{Inventory the bytes an agent workflow produces}
\fontsize{10.5}{12.5}\selectfont\setlength{\tabcolsep}{4pt}\renewcommand{\arraystretch}{1.15}
\begin{tabular}{>{\raggedright\arraybackslash}p{\dimexpr 0.23\linewidth-3.68pt\relax}>{\raggedright\arraybackslash}p{\dimexpr 0.77\linewidth-12.32pt\relax}}
\toprule
\textbf{Artifact} & \textbf{Concrete example to include in the map} \\ \midrule
Prompt + attachments & Geometry excerpts, material definitions, error traces and user annotations. \\[2pt]
Tool traffic & Search query text, remote shell arguments, embedding requests or uploaded files. \\[2pt]
Derived artifacts & Summaries, vector indexes, generated input decks and result tables. \\[2pt]
Operational copies & TUI session databases, server request logs, backups and support bundles. \\[2pt]
\bottomrule\end{tabular}\par\vspace{2mm}
\end{frame}
\begin{frame}[t,fragile]{Worked map: OpenCode plus a private LLM}
\fontsize{11.5}{14}\selectfont
\begin{itemize}\setlength{\itemsep}{2mm}
\item Laptop $\rightarrow$ inference: prompt, selected repository text, and previous tool results.
\item Laptop $\rightarrow$ solver: local files and a command; solver $\rightarrow$ laptop: logs and statepoint results.
\item Laptop $\rightarrow$ search: query terms, if search is enabled. This is an independent external transmission.
\item Record the allowed data classes for each arrow, the account involved, and the retained copies.
\end{itemize}
\end{frame}
\begin{frame}[t,fragile]{Ask providers questions that require evidence}
\fontsize{10.5}{12.5}\selectfont\setlength{\tabcolsep}{4pt}\renewcommand{\arraystretch}{1.15}
\begin{tabular}{>{\raggedright\arraybackslash}p{\dimexpr 0.2\linewidth-3.2pt\relax}>{\raggedright\arraybackslash}p{\dimexpr 0.8\linewidth-12.8pt\relax}}
\toprule
\textbf{Topic} & \textbf{A specific question} \\ \midrule
Training use & Are submitted prompts, files and outputs used for training under the applicable contract? \\[2pt]
Retention & Which request and operational records persist, for how long, and under what exceptions? \\[2pt]
Access & Which personnel, locations and subprocessors can access plaintext or recoverable copies? \\[2pt]
Isolation & Which technical and contractual controls separate our allocation and administration? \\[2pt]
Change control & How are changes to locations, access, terms or subprocessors communicated? \\[2pt]
\bottomrule\end{tabular}\par\vspace{2mm}
\end{frame}
\begin{frame}[t,fragile]{Encryption addresses a defined part of the path}
\fontsize{11.5}{14}\selectfont
\begin{itemize}\setlength{\itemsep}{2mm}
\item TLS protects transport between authenticated endpoints; it does not prevent the endpoint process from seeing plaintext.
\item Disk encryption protects stored media under its key-management assumptions; administrators may still have application-level access.
\item Confidential-computing designs require evidence about attestation, key release, GPU/CPU coverage, logs, and residual access.
\item Evaluate the threat model and legal authority separately; a product label is not the completed review.
\end{itemize}
\end{frame}
\begin{frame}[t,fragile]{Deployment options are information-flow controls}
\fontsize{10.5}{12.5}\selectfont\setlength{\tabcolsep}{4pt}\renewcommand{\arraystretch}{1.15}
\begin{tabular}{>{\raggedright\arraybackslash}p{\dimexpr 0.22\linewidth-5.28pt\relax}>{\raggedright\arraybackslash}p{\dimexpr 0.44\linewidth-10.56pt\relax}>{\raggedright\arraybackslash}p{\dimexpr 0.34\linewidth-8.16pt\relax}}
\toprule
\textbf{Option} & \textbf{Control you gain} & \textbf{What you give up} \\ \midrule
Commercial / cloud API & Provider terms decide retention, access and revision. & Fastest start; least control over where bytes travel. \\[2pt]
Facility-hosted & Endpoint, logs and access follow project policy. & The institution carries serving operations (module 06). \\[2pt]
On-premises / local & Bytes stay on approved hosts; the boundary is the machine. & You own the full stack: hardware, maintenance, recovery. \\[2pt]
\bottomrule\end{tabular}\par\vspace{2mm}
\end{frame}
\begin{frame}[t,fragile]{Write a technology-control proposal}
\fontsize{10.5}{12.5}\selectfont\setlength{\tabcolsep}{4pt}\renewcommand{\arraystretch}{1.15}
\begin{tabular}{>{\raggedright\arraybackslash}p{\dimexpr 0.2\linewidth-3.2pt\relax}>{\raggedright\arraybackslash}p{\dimexpr 0.8\linewidth-12.8pt\relax}}
\toprule
\textbf{Control} & \textbf{Concrete project implementation} \\ \midrule
Approved data & A manifest of permitted public/synthetic inputs plus owner and revision. \\[2pt]
Access & Named project accounts, least privilege, access review and removal procedure. \\[2pt]
Egress & Reviewed tool destinations; network controls where the project requires enforcement. \\[2pt]
Retention & Documented session/log storage, access, cleanup and backup handling. \\[2pt]
Change trigger & New material, collaborator, operator, model endpoint, or external tool. \\[2pt]
\bottomrule\end{tabular}\par\vspace{2mm}
\vfill{\fontsize{6.5}{8}\selectfont\color{UTgray} Sources: utexport. See references and instructor notes.\par}
\end{frame}
\begin{frame}[t,fragile]{Worked decision memo: what can be concluded?}
\fontsize{11.5}{14}\selectfont
\begin{itemize}\setlength{\itemsep}{2mm}
\item Known: the source report is public; the proposed endpoint and its stated retention terms are identified.
\item Unknown: whether a support subcontractor can access request logs and whether the private annotations are within an exclusion.
\item Recommendation: use only the public source for the class demonstration; hold the annotated case for institutional review.
\item Decision requested: determine the annotated material's status and whether the documented access arrangement is authorized.
\end{itemize}
\end{frame}
\begin{frame}[t,fragile]{Know what kind of answer a reviewer can give}
\fontsize{11.5}{14}\selectfont
\begin{itemize}\setlength{\itemsep}{2mm}
\item The institution can coordinate jurisdiction, project restrictions and the proposed handling arrangement.
\item Under \S{}810.5, DOE advice can be requested about scope and authorization; the provision distinguishes advice from a binding interpretation.
\item Retain the written response and the facts on which it relied.
\item Reopen the question when those facts change; do not treat a past answer as blanket approval for a different workflow.
\end{itemize}
\vfill{\fontsize{6.5}{8}\selectfont\color{UTgray} Sources: 810interpret, utexport. See references and instructor notes.\par}
\end{frame}
\begin{frame}[t,fragile]{If a workflow crosses the agreed boundary}
\fontsize{11.5}{14}\selectfont
\begin{itemize}\setlength{\itemsep}{2mm}
\item Stop the affected transfer or tool integration and preserve the factual timeline.
\item Record what material moved, to which destination, through which account, and which copies may remain.
\item Contact the institutional reporting/response channel and follow its instructions.
\item Do not let an agent delete evidence, invent a compliance conclusion, or contact third parties without an authorized response plan.
\end{itemize}
\end{frame}
\begin{frame}[t,fragile]{Lab: produce a reviewable boundary package}
\fontsize{11.5}{14}\selectfont
\begin{itemize}\setlength{\itemsep}{2mm}
\item Submit a data inventory, information-flow map, source/contract evidence, and a one-page decision memo.
\item Analyze a public-document case, a mixed public/private case, a private-LLM/external-tool case, and a multinational-team case.
\item For each, separate scope/exclusion analysis, authorization route, reporting ownership, and architecture controls.
\item Grade success by completeness and correct routing, including explicit unknowns and proposed interim controls.
\end{itemize}
\end{frame}
\begin{frame}[t,fragile]{Sample submission: the boundary package}
\fontsize{11.5}{14}\selectfont
\begin{itemize}\setlength{\itemsep}{2mm}
\item Package structure: a data inventory (one row per data class), an information-flow map, one section per case, a one-page decision memo; the cases are synthetic.
\item Public-document case: the scope analysis records the public status with source citations; the authorization route is stated as a working assumption pending institutional review.
\item Mixed public/private case: the private input determines the package status; the memo names the unknowns and proposes a synthetic substitute as the interim control.
\item Private-LLM/external-tool case: model-provider access is separated from tool-side destinations; controls include scoped tool permissions and a retained evidence manifest.
\item Multinational-team case: the access and recoverability questions drive the analysis; the interim control restricts the workflow to approved material.
\end{itemize}
\end{frame}
\begin{frame}[t]{References 1}
\fontsize{9}{12}\selectfont\begin{itemize}\setlength{\itemsep}{3mm}
\item \textbf{curriculum}: User-supplied PHANES curriculum: mod/01 through mod/08/README.md.
\item \textbf{colors}: \href{https://umac.utexas.edu/brand-center/colors/}{umac.utexas.edu/brand-center/colors}
\item \textbf{utexport}: \href{https://research.utexas.edu/resources/research-security/export-control}{research.utexas.edu/resources/research-security/export-control}
\item \textbf{810scope}: \href{https://www.ecfr.gov/current/title-10/chapter-III/part-810/section-810.2}{ecfr.gov/current/title-10/chapter-III/part-810/section-810.2}
\item \textbf{110}: \href{https://www.nrc.gov/reading-rm/doc-collections/cfr/part110/index}{nrc.gov/reading-rm/doc-collections/cfr/part110/index}
\item \textbf{earrelease}: \href{https://www.ecfr.gov/current/title-15/subtitle-B/chapter-VII/subchapter-C/part-734/section-734.13}{ecfr.gov/current/title-15/subtitle-B/chapter-VII/subchapter-C/part-734/section-734.13}
\item \textbf{810definitions}: \href{https://www.ecfr.gov/current/title-10/chapter-III/part-810/section-810.3}{ecfr.gov/current/title-10/chapter-III/part-810/section-810.3}
\end{itemize}\end{frame}
\begin{frame}[t]{References 2}
\fontsize{9}{12}\selectfont\begin{itemize}\setlength{\itemsep}{3mm}
\item \textbf{810general}: \href{https://www.ecfr.gov/current/title-10/chapter-III/part-810/section-810.6}{ecfr.gov/current/title-10/chapter-III/part-810/section-810.6}
\item \textbf{810specific}: \href{https://www.ecfr.gov/current/title-10/chapter-III/part-810/section-810.7}{ecfr.gov/current/title-10/chapter-III/part-810/section-810.7}
\item \textbf{810reports}: \href{https://www.ecfr.gov/current/title-10/chapter-III/part-810/section-810.12}{ecfr.gov/current/title-10/chapter-III/part-810/section-810.12}
\item \textbf{810interpret}: \href{https://www.ecfr.gov/current/title-10/chapter-III/part-810/section-810.5}{ecfr.gov/current/title-10/chapter-III/part-810/section-810.5}
\item \textbf{earresearch}: \href{https://www.ecfr.gov/current/title-15/subtitle-B/chapter-VII/subchapter-C/part-734/section-734.8}{ecfr.gov/current/title-15/subtitle-B/chapter-VII/subchapter-C/part-734/section-734.8}
\end{itemize}\end{frame}
\end{document}